\section{Coherent access: a tight bound for block encodings}\label{sec:coherent}

Quantum lower bounds for $b^*H^{-1}b$ depend on whether an algorithm can
query matrix entries or only a unitary with $H$ as a normalized block.
For the latter we prove a lower bound (Theorem~\ref{thm:coherent}) that
matches, up to logarithmic factors, the known upper bound of Chakraborty,
Gily\'en, and Jeffery, and we show that
the hard instances become easy with exact entry queries.

An $(\alpha,a,\delta_{\rm in})$ block encoding is a unitary $U_H$ such that
\[
 \norm{H-\alpha(\langle0^a|\otimes I)U_H(|0^a\rangle\otimes I)}
 \leq\delta_{\rm in}.
\]
The algorithm may call $U_H$, $U_H^\dagger$, and their controlled
versions.  A second controlled unitary prepares $|b\rangle=b/\norm b$, and
$\norm b$ is known.  In the \emph{plain block-encoding model}, these are
the only oracles.  The
normalization $\alpha\geq1$ is part of the input; it is not the
condition number.

\begin{theorem}[Relative inverse forms under block access]\label{thm:coherent}
Let $\kappa^{-1}I\preceq H\preceq I$, $\kappa\geq4$, and
$0<\epsilon\leq1/16$.  At constant failure probability, the worst-case
number of block queries needed to estimate $Q=b^*H^{-1}b$ to relative
error $\epsilon$ is
\[
 \Omega(\alpha\kappa/\epsilon)
 \quad\text{and}\quad
 \widetilde O(\alpha\kappa/\epsilon).
\]
The upper bound is a special case of the variable-time algorithm of
Chakraborty, Gily\'en, and Jeffery for negative-power norms
\citep[Theorem 33 in the full version]{ChakrabortyGilyenJeffery2019}.
With implementation costs $T_H,T_b$ for the two controlled
unitaries, its cost at failure probability $\zeta$ is
\begin{equation}\label{eq:cgj-cost}
 O\!\left(\frac{\log^3\kappa}{\epsilon}
 \left[\alpha\kappa(T_H+a)\log^2\frac\kappa\epsilon
       +\sqrt\kappa\,T_b\right]
 \log\frac{\log\kappa}{\zeta}\right).
\end{equation}
Their theorem applies here when
$\delta_{\rm in}=o(\epsilon/[\kappa\log^3(\kappa/\epsilon)])$.
The lower bound holds even for exact encodings in dimension three with
public $b$.
\end{theorem}

\begin{proof}
Set $R_b=\norm{H^{-1/2}|b\rangle}$.  The cited theorem with negative-power
parameter $c=1/2$ estimates $R_b$ multiplicatively.  If
$|\widehat R_b/R_b-1|\leq\epsilon/3$, then
$|\widehat R_b^2/R_b^2-1|\leq\epsilon$ for $\epsilon\leq1$.
Thus $\norm b^2\widehat R_b^2$ estimates $Q$, with the displayed cost.
This invocation estimates a norm directly; preparing the normalized
inverse state alone would not provide $Q$.

For the lower bound, let
\[
 H_z=\diag\left(\kappa^{-1},
           2(1+4\epsilon z)/\kappa,1\right),\qquad
 b=e_2,\qquad z\in\{0,1\}.
\]
Both matrices have condition number exactly $\kappa$, and
$Q_0=\kappa/2$, $Q_1=\kappa/[2(1+4\epsilon)]$.  Their relative-error
intervals are disjoint, since
$(1-\epsilon)(1+4\epsilon)>1+\epsilon$.
For a Hermitian contraction $A$, define the exact unitary completion
\[
 W(A)=\begin{pmatrix}A&\sqrt{I-A^2}\\
                    \sqrt{I-A^2}&-A\end{pmatrix}.
\]
Take $U_z=W(H_z/\alpha)$.  The only changing entry of $H_z/\alpha$
lies in $[0,5/8]$, where $x\mapsto\sqrt{1-x^2}$ has bounded
derivative.
Consequently
$\norm{U_0-U_1}=O(\epsilon/(\alpha\kappa))$.
Replacing the oracle calls one at a time bounds the Euclidean distance
between final purified states of any $T$-query algorithm by
$T\norm{U_0-U_1}$; inverse and controlled calls satisfy the same bound.
Measurements and adaptive operations can be deferred or purified.
Distinguishing the two disjoint output intervals with constant bias
requires constant trace distance, and hence $T=\Omega(\alpha\kappa/\epsilon)$.
\end{proof}

The lower bound fails for this family under exact sparse-value
access: one query to the second diagonal entry reveals $z$.  For coherent
$d$-sparse access, the standard block encoding has normalization at most
$d\norm H_{\max}\leq d$ \citep{GilyenEtAl2019}.  Together with
Corollary~\ref{cor:coherent-counting}, this bounds the worst-case number
$Q_{\rm sparse}$ of coherent $d$-sparse queries for relative-$\epsilon$
estimation of $Q$ over $d$-sparse $H$ with $\kappa^{-1}I\preceq H\preceq I$, with the preparation of $|b\rangle$ counted separately,
for $d\geq2$, $\kappa\geq4$, $0<\epsilon\leq1/128$, and dimensions
large enough for the counting family:
\begin{equation}\label{eq:sparse-quantum-gap}
 \Omega(\sqrt\kappa/\epsilon)
 \ \leq\ Q_{\rm sparse}\ \leq\
 \widetilde O(d\kappa/\epsilon).
\end{equation}
These bounds do not match in general.
The lower bound is proved for the two-sparse sign-block family, whose
matrices commute and have public eigenvalues; the hidden input changes
only how much weight $b$ places on each eigenvalue.  Its complexity
$\Theta(\sqrt\kappa/\epsilon)$ therefore does not
contradict Theorem~\ref{thm:coherent}, whose lower bound comes from an
unknown eigenvalue.

\subsection{Why a single polynomial transform loses a square-root factor}

A simpler, weaker algorithm block-encodes
$B=H^{-1/2}/(2\sqrt\kappa)$ and applies it to $|b\rangle$.  Its
success probability is
\[
 p_b=\frac{\langle b|H^{-1}|b\rangle}{4\kappa}
 \in[1/(4\kappa),1/4].
\]
Amplitude estimation \citep{BrassardEtAl2002} gives additive error
$O(\sqrt{p_b}/L+L^{-2})$ after $L$ uses of this circuit.  Relative error
$\epsilon$ therefore needs $O(\sqrt\kappa/\epsilon)$ uses in the worst
case.  An error $\xi=O(\epsilon/\sqrt\kappa)$ in the block encoding of $B$
suffices, because it changes the success probability by at most
$2\xi\sqrt{p_b}+\xi^2$.  A negative-power polynomial approximation
bounded on the whole signal domain $[-1,1]$ \citep{GilyenEtAl2019}
implements this transform with $\widetilde O(\alpha\kappa)$ block queries,
for a total cost of $\widetilde O(\alpha\kappa^{3/2}/\epsilon)$.
The variable-time algorithm behind \eqref{eq:cgj-cost} avoids
multiplying these two worst-case costs.

\begin{proposition}[Bounded-transform obstruction]\label{prop:bounded-transform}
Suppose $\kappa\geq8$, $\alpha\geq1$, and a degree-$n$ complex polynomial
$P$ satisfies $|P(x)|\leq1$ for every $x\in[-1,1]$.  If $P$ approximates
$1/(2\sqrt{\alpha\kappa x})$ to absolute error $1/32$ at both
$x_1=1/(\alpha\kappa)$ and $x_2=4/(\alpha\kappa)$, then
$n=\Omega(\alpha\kappa)$.
The same conclusion holds for a bounded Laurent transform of the signal
phase with bandwidth $n$.
\end{proposition}

\begin{proof}
The target values are $1/2$ and $1/4$, so the real part of $P$ changes
by at least $3/16$ on an interval of length $3/(\alpha\kappa)$.
At some $\xi\in[x_1,x_2]\subset[0,1/2]$, its derivative has magnitude
at least $\alpha\kappa/16$.  Bernstein's polynomial inequality gives
$|\operatorname{Re}P'(\xi)|\leq n/\sqrt{1-\xi^2}\leq2n/\sqrt3$.
For a Laurent transform $F(e^{i\theta})$ bounded by one on the unit
circle, the two phases $\arccos x_1,\arccos x_2$ are separated by
$\Theta(1/(\alpha\kappa))$.  The trigonometric Bernstein inequality
$\norm{dF/d\theta}_\infty\leq n\norm F_\infty$ gives the same bound.
\end{proof}

The proof uses the elementary Bernstein inequalities only for bounded
polynomials and trigonometric polynomials.  The obstruction is interior: the relevant eigenvalues of $H/\alpha$
occupy a small part of the domain on which the transform must stay
bounded.  Removing the parity restriction of QSVT, for example by
generalized signal processing \citep{MotlaghWiebe2024}, keeps this
boundedness and so does not remove the obstruction.  The proposition gives no lower bound for variable-time,
rational, or postselected algorithms.

Orsucci and Dunjko also show that the access model matters
\citep[Proposition 6 and Sections 4.2--4.3]{OrsucciDunjko2021}.
They prove a worst-case $\Omega(\kappa)$ query lower bound for preparing
solution states of positive definite systems under their matrix oracles.
With a normalized encoding of $I-\eta A$ instead, their inverse
approximation has degree $\widetilde O(\sqrt\kappa)$ for
fixed $\eta>0$, because its steep part lies near an
endpoint of the signal domain.  This shifted encoding is an additional
access assumption; obtaining it from an arbitrary
encoding of $A$ is not free.  Their output is a solution state and their construction
uses shifted access, whereas we estimate a scalar to relative error with
an unshifted bounded transform.  The square-root polynomial degree
also omits their state-preparation cost.

\subsection{Relation to earlier scalar quantum algorithms}

Scalar readout of a linear-system solution is prior work.  Alase et al.\ \citep{AlaseEtAl2022} study
additive estimation of $x^*Cx$ for $Ax=b$.  For expectation values in a
given state they prove
the tight complexity $\Theta(\alpha_C/\epsilon)$ in queries to an encoding
of the \emph{observable} $C$, and use it to lower-bound the
linear-system scalar task.  They also explicitly distinguish block and entry access
for a fixed computational-basis state: their block-query lower
bound still holds, whereas one exact-entry query suffices
\citep[discussion after Theorem IV.13 and Appendix A]{AlaseEtAl2022}.
Our lower bound counts queries to the inverted matrix $H$, concerns
relative error, and gains an additional factor $\kappa$ from an
eigenvalue that depends on the hidden bit.  We claim neither the first
scalar-output quantum algorithm nor the negative-power estimator.
Theorem~\ref{thm:coherent} contributes the matching condition-number
dependence for relative estimation of the inverse form, and an
explicit example of the block--entry access difference for
this task.  Likewise, a constant-query quantum decision of an underlying
Forrelation promise is not automatically a constant-query estimator of
every associated inverse form; see the discussion of $q_*$ after
Proposition~\ref{prop:rational-clock}.
